% !PRINTED 453
\weberpage{453}

%<<CONT
man die $P_8$ passend die \gesperrt{Hauptpole} nennen und die durch sie gehenden Axen mit $A_8, A_6, A_4, A_2$ bezeichnen.

Hervorzuheben sind ferner die 28 Paare von Verbindungslinien eines $P_6$ mit den beiden zugehörigen $P_3$, die wir mit $A_3$ bezeichnen können, und endlich die 24 Verbindungslinien je zweier zusammengehöriger $P_7$, als deren Repräsentanten die Seiten des Coordinatendreiecks zu betrachten sind, die mit $A_7$ bezeichnet werden können.

\begin{center}
\S.\,121.\\
\gesperrt{Invariantencurven der Gruppe $G_{168}$.}
\end{center}

Eine Form $\mu^{\mathrm{ten}}$ Grades $\varphi\,(x_1, x_2, x_3)$ wird durch die Substitutionen der Gruppe im Allgemeinen in 168 verschiedene Formen übergehen. Bei besonderen Formen von $\varphi$ kann diese Zahl sich aber verringern, und dann bilden die Substitutionen, durch die $\varphi$ ungeändert bleibt, eine in $G_{168}$ enthaltene Gruppe $G'$, deren Grad ein Theiler von 168 sein muss. Die Anzahl der Formen, in die $\varphi$ übergeht, ist der Index des Theilers $G'$, und jede dieser verschiedenen Formen von $\varphi$ bleibt durch eine mit $G'$ conjugirte Gruppe ungeändert. Die Form $\varphi$ ist eine absolute Invariante der Gruppe $G'$.

Die Gleichung $\varphi = 0$ stellt eine auf das Coordinatendreieck $x_1, x_2, x_3$ bezogene Curve dar, und diese Curve wird eben durch die Substitutionen von $G_{168}$ auf andere Curven abgebildet. Sind die 168 Bildcurven nicht alle von einander verschieden, so bleibt die Curve $\varphi$ durch die Substitutionen einer Gruppe $G'$ ungeändert oder ändert sich nur um einen constanten Factor, ist also relative Invariante von $G'$. Eine gerade Linie bleibt nur dann durch andere als die identische Substitution ungeändert, wenn sie zu den im vorigen Paragraphen beschriebenen Axen gehört.

Alle Eigenschaften und Beziehungen zwischen Punkten, Linien und Curven, die durch lineare Transformation unzerstörbar sind (die sogenannten projectiven Eigenschaften der Geometrie), bleiben in den Bildern erhalten. Wenn also z.\,B. eine gerade Linie Tangente oder Wendetangente oder Doppeltangente einer Curve ist, so stehen alle Bilder der geraden Linie in derselben Beziehung zu den Bildern der Curve. Ebenso wenn ein Punkt
%CONT>>
